\problemname{Segmentövervakning}

Efter att ha lyckats erövra alla stadsstater i Kattago, vill det övermänskliga kattriket, Katen, och dess ledare, 
Kattis, ha bättre kontroll över sina erövrade stadsstaterna, men specifikt gatorna.
För att undivka kat(t)astrof, har Kattis bett företaget Purrfect Segments Inc. att producera olika kameror som 
kan övervaka olika segment på en gata. 

Företaget fick uppdraget av Kattis att designa och installera kameror på en gata i Katen som kan representeras som en tallinje. 
Purrfect Segments Inc. har utvecklat sina kamerors paw-sition så att kamerorna kan specifikt övervaka ett segment $(a,b)$ till purr-fektion, 
vilket innebär att kameran övervakar alla heltal $x$ på tallinjen sådana att $a<=x<=b$. 
Trots företagets välutvecklade kameror, så har de ändå stött på flertal problem, bland annat att många hårdvaror i kamerorna behöver uppdateras regelbundet, 
vilket innebär att installerade kameror kommer möjligtvis behöver tas ner efter ett antal dagar.

Under $Q$ dagar, kommer företaget antingen installera en kamera med serienummer $i$ som övervakar ett segment $(a,b)$ på gatan, 
eller ta bort en specifik kamera från gatan som innan var installerat. 

Men under vissa dagar är det väldigt onödigt att aktivera alla installerade kameror, 
och det är mer miljövänligt att endast aktivera ett fåtal kameror. 

Kattis vill inte %TODO ändra statement så att den passar in ändringen

Som tur är har Purrfect Segments Inc. lyckats tänka lite i förväg, och kan alltid garantera att man endast behöver aktivera 
maximalt 2 kameror varje dag, så att så många heltal övervakas som möjligt. 
Alltså för en specifik dag, kallar vi det minsta talet som övervakas av någon nuvarande installerad kamera för $x_{min}$, 
och vi kallar det största talet som övervakas av någon nuvarande installerad kamera för $x_{max}$. 
Då är det garanterat att alla tal mellan $x_{min}$ och $x_{max}$ (inklusive $x_{min}$ och $x_{max}$) kan bli övervakade genom att endast aktivera max 2 kameror.

Från början finns det inga kameror som övervakar denna gata i Katen. För varje dag som går upp till den $Q:te$ dagen:
Vilka kameror ska företaget aktivera för att maximera sträckan som är övervakad varje dag?

\begin{figure}[h]
  \centering
\includegraphics[width=0.25\textwidth]{kattIsWatching.png}
  \caption{Figur till vänster associeras vanligtvis med "kiki". Figur till höger associeras vanligtvis med "boba".}
  \label{fig:kikiboba}
\end{figure}


\section*{Indata}
Den första raden i indatan innehåller heltalet $Q$ ($1 \le Q \le 2 \cdot 10^5$), antal dagar i testfallet.

De följande $Q$ rader börjar antingen med ``+'' ifall en kamera ska installeras, eller ``-'' ifall en kamera ska tas bort:
\being{itemize}
  \item ``+'': Raden följer med 3 heltal, där det första heltalet består av serienumret på kameran $s$ $(1\le s \le 10^6)$, 
  och det andra och tredje heltalet, a och b $(0\le a \le b \le 10^9)$, är start respektive slut på det segment som kameran övervakar.
  \item ``-'': Raden följer med 1 heltal, där heltalet är serienumret på kameran $s$ som ska tas bort från gatan. Det är garanterat att kameran som tas bort är installerad.
\end{itemize}

Alla kameror har ett unikt serienummer. Det är även garanterat att för varje dag finns det alltid minst 1 kamera som är installerad på gatan.

\section*{Utdata}
Skriv ut $Q$ rader, där det $i$:te raden är svaret till den $i$:te dagen: 
Det minsta antalet kameror som behöver aktiveras den dagen. Antingen 1, 2, eller -1 ifall det är omöjligt att uppfylla reglerna 

\section*{Poängsättning}
Din lösning kommer att testas på en mängd testfallsgrupper.
För att få poäng för en grupp så måste du klara alla testfall i gruppen.

\noindent
\begin{tabular}{| l | l | p{12cm} |}
  \hline
  \textbf{Grupp} & \textbf{Poäng} & \textbf{Gränser} \\ \hline
  $1$    & $20$       & $Q \le 2$ \\ \hline
  $2$    & $15$       & Svaret är antingen 1 eller ``impossible''. \\ \hline
  $3$    & $30$       & $Q \le 1000$. \\ \hline
  $4$    & $35$       & Inga ytterligare begränsningar. \\ \hline
\end{tabular}

\section*{Förklaring av exempelfall 1}
Ordet innehåller 2 st ``\texttt{b}'', medan det finns inga ``\texttt{k}''. Det finns alltså fler ``\texttt{b}'' än ``\texttt{k}'', och svaret är därför ``\texttt{boba}''.

\section*{Förklaring av exempelfall 3}
Ordet innehåller 1 st ``\texttt{b}'' och 1 st ``\texttt{k}''. Det finns lika många ``\texttt{b}'' som ``\texttt{k}'', och därför är svaret ``\texttt{boki}''.
